\section{Full-jet transfer for the same source--relation volume}
\label{sec:confluent-source-transfer}

\textbf{Purpose and result.}
The source and relation determinants in the preceding chapters can be
computed from a fixed-width family of complete derivative jets. This
chapter gives the full reviewed construction, its divisibility proof,
the positive relation norms, the phase solve and the exact maps from the
auxiliary squared relation to the actual second conormal quotient.
It does not identify those two quotients: the displayed surjections and
their kernels calculate their relation.

For the original dagger-stable polynomial, with $c=k/2$ and
$\psi(u)=i^{-q}\chi(c+iu)$, the coordinate isomorphism is
\[
\Psi:P(S)\longmapsto P(c+iu),\qquad
\Psi(\chi^2)=i^{2q}\psi^2,\qquad
(\partial_u^d\Psi P)(\zeta_a)=i^dP^{(d)}(\lambda_a).
\]
No derivative factorial, multiplicity, unit or phase has been dropped.
With the source norms $\omega_j$, relation norms $\nu_j$ and full jet
matrix $F_n$, the exact determinant calculation is
\[
\frac{\mathfrak B_n}{\mathfrak D_n}
=\frac{\det F_n}{\det V},\qquad
\det G_N=\left(\prod_{j=N-q+1}^{N}\omega_j\right)
\frac{\det V}{\det F_{N-q+1}},\quad N\ge q-1.
\]
The raw confluent matrix $V$ below is not the quotient-volume scalar $V_N$.
The input remains the original arithmetic convolution measure.

\textbf{Source and checking scope.}
The complete corrected note below is the PR~24 text at
\href{https://github.com/KokunoYumeto/zeta-function-research-reader/tree/b8d0e04e06854c2e80b288041889659ab2267ebe/workbenches/tau-confluent-transfer}{its reviewed corrected revision}.
It was merged at \path{5e67416c467166cb8729f47cea941c22a9f5f33c}.
Its opening main/PR inspection statements and its expanded conversation
validation summary retain their historical meanings. Fresh integration
verification fetched the corrected public core and passed eight exact
methods in ordinary and optimized Python; both deliberate-failure
controls failed. No triggered CI was observed for the corrected head or
merge, and no new Lean or arithmetic interval certificate is claimed.
The formula blocks retain the source's literal coordinate notation,
with line wrapping only for readability.

\begingroup
\setlength{\emergencystretch}{3em}
\providecommand{\tightlist}{\setlength{\itemsep}{0pt}\setlength{\parskip}{0pt}}
\subsection[\hspace{0.4em}Scope and current-main
integration]{Scope and current-main
integration}\label{confluent-transfer-scope-and-current-main-integration}

This add-only analytic continuation starts from main
\texttt{fa4be32f\allowbreak{}87a9a90f\allowbreak{}3c02a07f\allowbreak{}7790dea9\allowbreak{}5a1e7b0f}, after the merge of PR
\#21 and its workflow-trigger correction. The PR \#22 mathematical note
was read at \texttt{811210d2\allowbreak{}4b80813a\allowbreak{}08972ca2\allowbreak{}3f919db0\allowbreak{}15383137}; its
strict verification was in progress at that reading. This contribution
neither modifies that branch nor treats pending checks as a certificate.

The source is the delivered Toda-volume continuation, following the
cyclic and exterior-trace notes. Its full LaTeX was read directly. The
new result is a fixed-width, full-jet transfer for the SAME canonical
source/relation quotient metric. It applies the classical Christoffel
transformation; it is not a claim to a new general orthogonal-polynomial
theorem or to RH.

Keep G(Z)-\textgreater Z, with its infinite target, and its coefficient
square to G(C)-\textgreater C. The structural base remains b\_tau. Every
supported coefficient map below retains e=0\^{}bullet and external tau
separately. The original theta complex is C\_+={[}V
-\/-Theta-\/-\textgreater{} B{]}, with D=-x d/dx, Q=B/Theta V and g=2
xi. No scalar, source function or measure is assigned a new
normalization.

For a nonempty finite packet h of actual zeros with full orders, retain
F\_h with Mellin transform g/h. For the doubled-original-root transfer
below, require the same dagger stability as the original source: its
cyclic annihilator satisfies the polynomial identity
conj(chi)(k-S)=(-1)\^{}q chi(S). The complete reflection-stable quartet
packets used in the final target satisfy this hypothesis. It does not
follow from being an arbitrary finite subpacket of zeros. At tensor
degree k\textgreater=1, keep the actual cyclic sum module

\begin{Verbatim}[breaklines=true,breakanywhere=true,fontsize=\small]
C=C[S]/(chi),  q=deg chi>=1,  c=k/2,  A=M_S,
eta[P]=upsilon_h^(tensor k) P(sum_i M_(s_i))1,
upsilon_h=j_h(g/h),
Vcal P=P(D_1+...+D_k)(F_h^(tensor k)).
\end{Verbatim}

The original equations are

\begin{Verbatim}[breaklines=true,breakanywhere=true,fontsize=\small]
J^(k) Vcal=eta pi_chi,
q^(k) Vcal=sigma_h^(tensor k) eta pi_chi.
\end{Verbatim}

Its norm is the actual integral against m\_k=w\_h\^{}\{*k\}, where

\begin{Verbatim}[breaklines=true,breakanywhere=true,fontsize=\small]
w_h(t)=|(g/h)(1/2+i t)|^2/(2 pi),  integral m_k=mu_h^k.
\end{Verbatim}

The auxiliary measure is dmu\_theta(u)=exp(theta u)m\_k(u)du,
\textbar theta\textbar\textless pi/2. Its mass is M\_h(theta)\^{}k.
Theta in this expression is the real tilt parameter, not the
theta-summation map. The exponential-moment bound from the source makes
all finite derivatives below valid. At tilt zero the metric is the
original arithmetic metric.

For h=1 the finite arithmetic quotient is zero. Its determinant is the
empty determinant 1; the analytic seed and its mass remain, and G(0)
still has e and tau. The positive-dimensional inverses below require
q\textgreater=1.

\subsection[\hspace{0.4em}Exact coordinate and raw-jet
maps]{Exact coordinate and raw-jet
maps}\label{confluent-transfer-exact-coordinate-and-raw-jet-maps}

Use the ring isomorphism

\begin{Verbatim}[breaklines=true,breakanywhere=true,fontsize=\small]
Psi:C[S]->C[u],  P |-> P(c+i u),  u |-> (S-c)/i.
\end{Verbatim}

Thus d\_u Psi(P)=i Psi(P\textquotesingle), and multiplication by S
remains c+i M\_u. Define the original norm polynomial

\begin{Verbatim}[breaklines=true,breakanywhere=true,fontsize=\small]
Pi(u)=conj(chi)(c-i u) chi(c+i u).
\end{Verbatim}

Its leading coefficient is literally (-i)\^{}q i\^{}q=1. It is real and
nonnegative on R. Reflection gives chi\^{}dagger=(-1)\^{}q chi, hence

\begin{Verbatim}[breaklines=true,breakanywhere=true,fontsize=\small]
psi(u)=i^(-q) chi(c+i u) is real and monic,
Pi=psi^2,  Psi(chi)=i^q psi.
\end{Verbatim}

The factor i\^{}q remains in the source multiplication map. For distinct
roots lambda\_a of chi, of complete lengths ell\_a, put

\begin{Verbatim}[breaklines=true,breakanywhere=true,fontsize=\small]
zeta_a=(lambda_a-c)/i,  r_a=2 ell_a,  r=sum r_a=2q.
\end{Verbatim}

Explicitly, the coordinate isomorphism Psi induces C{[}S{]}/chi\^{}2
-\textgreater{} C{[}u{]}/Pi because Psi(chi\^{}2)=\allowbreak{}i\^{}(2q)Pi, and the
raw jet conversion is (d\_u\^{}d Psi(P))(zeta\_a)=i\^{}d
P\^{}(d)(lambda\_a). This is the actual map used below, including every
derivative factor. Without the displayed dagger identity, Pi need not
equal psi\^{}2 and its roots need not be the original zeta\_a with
doubled orders; that unsupported general-packet version is not asserted
here.

Define raw, NOT factorial-divided, derivative coordinates

\begin{Verbatim}[breaklines=true,breakanywhere=true,fontsize=\small]
J_2 f=(f^(d)(zeta_a))_(a,0<=d<r_a).
\end{Verbatim}

Their kernel is (Pi). Their product uses the full binomial rule.
Multiplication by u has matrix

\begin{Verbatim}[breaklines=true,breakanywhere=true,fontsize=\small]
(X_Pi v)_(a,d)=zeta_a v_(a,d)+d v_(a,d-1),  v_(a,-1)=0.
\end{Verbatim}

All nilpotent orders and the coefficient d are retained.

Let Q\_j be the real monic orthogonal source polynomials for dmu\_theta,
obtained by subtracting lower-degree projections, and write

\begin{Verbatim}[breaklines=true,breakanywhere=true,fontsize=\small]
omega_j=integral Q_j^2 dmu_theta>0,
a_j=omega_j/omega_(j-1),  a_0=0,
b_j=integral u Q_j^2 dmu_theta / omega_j.
\end{Verbatim}

The source recurrence is

\begin{Verbatim}[breaklines=true,breakanywhere=true,fontsize=\small]
u Q_j=Q_(j+1)+b_j Q_j+a_j Q_(j-1),  Q_(-1)=0.
\end{Verbatim}

No Q\_j is divided by its norm. Set

\begin{Verbatim}[breaklines=true,breakanywhere=true,fontsize=\small]
v_j=J_2 Q_j,
F_n=[v_n,...,v_(n+r-1)],
V=[J_2 1,...,J_2 u^(r-1)].
\end{Verbatim}

Both are r by r. The raw confluent determinant is

\begin{Verbatim}[breaklines=true,breakanywhere=true,fontsize=\small]
det V=(product_a product_(d=0)^(r_a-1) d!)
       product_(a<b)(zeta_b-zeta_a)^(r_a r_b).
\end{Verbatim}

The fixed row order is part of this formula. F\_0=V times the actual
unit-triangular polynomial coefficient matrix; only their determinants
are equal.

\subsection[\hspace{0.4em}Invertibility and the original relation
polynomial]{Invertibility and the original relation
polynomial}\label{confluent-transfer-invertibility-and-the-original-relation-polynomial}

\textbf{Invertibility proof.} Suppose F lies in
span(Q\_n,...,Q\_(n+r-1)) and J\_2 F=0. Then F=Pi H with deg
H\textless n. Orthogonality gives

\begin{Verbatim}[breaklines=true,breakanywhere=true,fontsize=\small]
0=<H,F>=integral Pi |H|^2 dmu_theta.
\end{Verbatim}

The integrand is nonnegative and its weight is positive almost
everywhere except finitely many points. Thus H=0, then F=0. For n=0 the
degree argument suffices. Hence every F\_n is invertible. This does NOT
estimate its inverse uniformly.

Solve

\begin{Verbatim}[breaklines=true,breakanywhere=true,fontsize=\small]
d_n=F_n^(-1)v_(n+r),   vartheta_n=-(d_n)_0.
\end{Verbatim}

The polynomial

\begin{Verbatim}[breaklines=true,breakanywhere=true,fontsize=\small]
Qhat_n=(Q_(n+r)-sum_(j=0)^(r-1)(d_n)_j Q_(n+j))/Pi
\end{Verbatim}

exists by exact raw-jet divisibility and is monic of degree n. Pairing
its numerator with every polynomial of degree below n proves that
Qhat\_n is precisely the monic orthogonal polynomial for the ORIGINAL
relation weight Pi dmu\_theta. Its norm is

\begin{Verbatim}[breaklines=true,breakanywhere=true,fontsize=\small]
nu_n=integral Qhat_n^2 Pi dmu_theta
    =-(d_n)_0 omega_n=omega_n vartheta_n>0.               (1)
\end{Verbatim}

Indeed only the Q\_n column pairs with the degree-n quotient; its
coefficient is -(d\_n)\_0. This proves vartheta\_n\textgreater0 with
arbitrary nonreal and repeated nodes.

Let K\_n={[}e\_1,...,e\_(r-1),d\_n{]}. Then

\begin{Verbatim}[breaklines=true,breakanywhere=true,fontsize=\small]
F_(n+1)=F_n K_n,
det K_n=(-1)^(r-1)(d_n)_0=vartheta_n.                  (2)
\end{Verbatim}

The sign uses the displayed even r=2q. On the same jet algebra,

\begin{Verbatim}[breaklines=true,breakanywhere=true,fontsize=\small]
T_n=F_n^(-1)X_Pi F_n,
T_(n+1)=K_n^(-1)T_n K_n,
d_n=(T_n-b_(n+r-1)I)e_(r-1)-a_(n+r-1)e_(r-2).         (3)
\end{Verbatim}

Thus the full operator and its jets have an explicit discrete transport,
rather than only a determinant observation.

\subsection[\hspace{0.4em}Fixed-width determinant and quotient-volume
formulas]{Fixed-width determinant and quotient-volume
formulas}\label{confluent-transfer-fixed-width-determinant-and-quotient-volume-formulas}

Let D\_n be the n by n source Hankel determinant, and B\_n the
corresponding determinant for Pi dmu\_theta. D\_0=B\_0=1. Orthogonality
gives

\begin{Verbatim}[breaklines=true,breakanywhere=true,fontsize=\small]
D_n=product_(j<n) omega_j,  B_n=product_(j<n) nu_j.
\end{Verbatim}

Multiplying (1) and (2) proves the complete confluent Christoffel
identity

\begin{Verbatim}[breaklines=true,breakanywhere=true,fontsize=\small]
B_n/D_n=product_(j<n) vartheta_j=det F_n/det V.          (4)
\end{Verbatim}

This is a proof from the source division and norm, not a distinct-root
approximation. Raw derivative factorials remain in det V.

For N\textgreater=q-1 set n=N-q+1. The source sequence is still

\begin{Verbatim}[breaklines=true,breakanywhere=true,fontsize=\small]
0 -> P_(N-q) --chi--> P_N -> C[S]/chi -> 0.
\end{Verbatim}

Its canonical least-norm quotient Gram G\_N satisfies the existing Schur
identity det G\_N=D\_(N+1)/B\_n. Consequently

\begin{Verbatim}[breaklines=true,breakanywhere=true,fontsize=\small]
V_N:=det G_N=(product_(j=n)^(n+q-1) omega_j) det V/det F_n. (5)
\end{Verbatim}

The constant coefficient map S=c+i u has determinant of modulus one on
each remainder basis, so these are the original quotient determinants,
with that coordinate map retained.

For n\textgreater=1,

\begin{Verbatim}[breaklines=true,breakanywhere=true,fontsize=\small]
delta_N=V_N/V_(N-1)=omega_N/nu_(n-1)
       =omega_N/(omega_(n-1) vartheta_(n-1)),          (6)
R_N:=V_(N-1)/V_(N+1)
   =omega_(n-1)omega_n vartheta_(n-1)vartheta_n
      /(omega_N omega_(N+1)).                         (7)
\end{Verbatim}

The transfer width 2q does not grow with N. The source recurrence and
high-degree polynomial evaluation are still required; no
computational-complexity or conditioning claim is inferred from the
width alone.

\subsection[\hspace{0.4em}One additional solve retains the full
phase]{One additional solve retains the full
phase}\label{confluent-transfer-one-additional-solve-retains-the-full-phase}

Differentiating actual source orthogonality under exp(theta u) gives

\begin{Verbatim}[breaklines=true,breakanywhere=true,fontsize=\small]
d_theta Q_j=-a_j Q_(j-1),  d_theta log omega_j=b_j.
\end{Verbatim}

For the first identity, every coefficient below j-1 vanishes by
differentiated orthogonality; pairing with Q\_(j-1) gives -omega\_j. Let
L\_n have entries (L\_n)\emph{(j-1,j)=-a}(n+j) for
1\textless=j\textless r, all other entries zero. Then

\begin{Verbatim}[breaklines=true,breakanywhere=true,fontsize=\small]
F_n'=F_n L_n-a_n v_(n-1)e_0^T.                        (8)
\end{Verbatim}

At n=0 the last term is zero. Trace of L\_n is zero. Hence

\begin{Verbatim}[breaklines=true,breakanywhere=true,fontsize=\small]
(log det F_n)'=-a_n e_0^T F_n^(-1)v_(n-1),
(log V_N)'=sum_(j=n)^(n+q-1) b_j
             +a_n e_0^T F_n^(-1)v_(n-1).             (9)
\end{Verbatim}

Retain T=(A-cI)/i and sigma=Tr T, which is real. The previous source
cross term is exactly

\begin{Verbatim}[breaklines=true,breakanywhere=true,fontsize=\small]
(b_N^*G_N b_(N+1))/omega_N=i(sigma-(log V_N)').
\end{Verbatim}

Here b\_N in that formula denotes the remainder vector of the S-monic
source polynomial, not the scalar recurrence b\_j. The factor i and the
distinction of types remain. The phase is phi\_N=sigma-(log
V\_N)\textquotesingle. At zero tilt for a full quartet, evenness gives
phi\_N=0; it is not set to zero at a general tilt.

\subsection[\hspace{0.4em}The positive gaps are norms of actual source
vectors]{The positive gaps are norms of actual source
vectors}\label{confluent-transfer-the-positive-gaps-are-norms-of-actual-source-vectors}

Put

\begin{Verbatim}[breaklines=true,breakanywhere=true,fontsize=\small]
Z_j=psi Qhat_j-Q_(j+q) in P_(j+q-1).
\end{Verbatim}

The leading terms cancel. Source orthogonality therefore proves

\begin{Verbatim}[breaklines=true,breakanywhere=true,fontsize=\small]
Delta_j:=||Z_j||^2=nu_j-omega_(j+q)>=0.                (10)
\end{Verbatim}

The first term psi Qhat\_j is, through the explicit unit i\^{}(-q), an
original chi relation. Its quotient is its supported zero. The
difference Z\_j instead has class -{[}Q\_(j+q){]} in the original
quotient. Its norm is not assigned to an absent element.

The SAME canonical rank-two control is W\_N=A*G\_N+G\_N A-kG\_N.
Combining its previously proved radius identity with (6) gives, for
n\textgreater=1,

\begin{Verbatim}[breaklines=true,breakanywhere=true,fontsize=\small]
epsilon_N^2=Delta_(n-1)Delta_n/(nu_(n-1)omega_N)-phi_N^2. (11)
\end{Verbatim}

Thus every factor in the upper bound

\begin{Verbatim}[breaklines=true,breakanywhere=true,fontsize=\small]
epsilon_N <= sqrt(Delta_(n-1)Delta_n/(nu_(n-1)omega_N))
\end{Verbatim}

comes from a specified positive source norm. Both gaps are computed by
the transfer:

\begin{Verbatim}[breaklines=true,breakanywhere=true,fontsize=\small]
Delta_(n-1)=omega_(n-1)vartheta_(n-1)-omega_N,
Delta_n=omega_n vartheta_n-omega_(N+1).
\end{Verbatim}

The previous two-step bound is now

\begin{Verbatim}[breaklines=true,breakanywhere=true,fontsize=\small]
epsilon_N <= sqrt(omega_(N+1)/omega_N)
                  (R_N-1)/(2 sqrt(R_N)),              (12)
\end{Verbatim}

with R\_N from (7). At the first degree N=q-1, use the separate
expression

\begin{Verbatim}[breaklines=true,breakanywhere=true,fontsize=\small]
epsilon_(q-1)^2=(nu_0-omega_q)/omega_(q-1)-phi_(q-1)^2.
\end{Verbatim}

No inverse at the previous rank-deficient degree is introduced.

\subsection[\hspace{0.4em}Keep the actual conormal quotient, not a repeated first
ideal]{Keep the actual conormal quotient, not a repeated first
ideal}\label{confluent-transfer-keep-the-actual-conormal-quotient-not-a-repeated-first-ideal}

The helper algebra for the norm computation is C{[}u{]}/Pi, isomorphic
through S=c+i u to C{[}S{]}/chi\_1\^{}2. The actual depth-two object is
C{[}S{]}/chi\_2, defined using the literal pullback of the original
I\^{}2. The merged conormal work gives chi\_2 \textbar{} chi\_1\^{}2.
The actual maps are

\begin{Verbatim}[breaklines=true,breakanywhere=true,fontsize=\small]
C[S]/chi_1^2 ->> C[S]/chi_2 ->> C[S]/chi_1.            (13)
\end{Verbatim}

Their kernels are respectively (chi\_2)/(chi\_1\^{}2) and
(chi\_1)/(chi\_2). In raw jets they truncate at the indicated full
orders and intertwine multiplication by S. Differentiation into depth
one also commutes, since chi\_1 \textbar{} chi\_2\textquotesingle. The
coordinate relation d\_u Psi=i Psi d\_S is retained.

For h=(s-rho)\^{}2 and tensor degree two, X=S-2rho gives

\begin{Verbatim}[breaklines=true,breakanywhere=true,fontsize=\small]
C[X]/X^6 ->> C[X]/X^5 ->> C[X]/X^3.
\end{Verbatim}

The first kernel is the class of X\^{}5. The retained depth-two class
X\^{}3 has derivative 3X\^{}2 at depth one. These maps use one original
scalar support. Their split lifts send quotient-killed classes to e at
the receiving label and fix external tau.

The full multiplier U=product\_i(g/h)(s\_i) retains the checked equation
delta(UP)=U delta(P)+(delta U)P. The second term is not discarded or
presumed cyclic. The 2q-jet norm calculation is not counted as
construction of new actual spectral multiplicities.

\subsection[\hspace{0.4em}Residual-controlled numerical
interface]{Residual-controlled numerical
interface}\label{confluent-transfer-residual-controlled-numerical-interface}

For an actual F and a specified candidate inverse Y, suppose a justified
bound gives
\textbar\textbar I-YF\textbar\textbar\_infinity\textless=eta\textless1.
For candidate d\_hat and residual r=v-F d\_hat, the exact Neumann
argument proves

\begin{Verbatim}[breaklines=true,breakanywhere=true,fontsize=\small]
||d-d_hat||_infinity <= ||Y||_infinity ||r||_infinity/(1-eta). (14)
\end{Verbatim}

This bounds the first coordinates used in (1) and (9). Arithmetic
application requires certified input enclosures for the moments, source
polynomials and zeros. Floating-point output alone is not a certificate.

Positive enclosures of omega and vartheta bound (7). Since
(R-1)\^{}2/(4R) is increasing for R\textgreater=1, (12) yields a valid
upper enclosure of epsilon\^{}2. The full phase can improve it through
(11). Finite positivity/invertibility is not substituted for a uniform
estimate of these residuals or inverse norms.

\subsection[\hspace{0.4em}Exact calibration and analytic
target]{Exact calibration and analytic
target}\label{confluent-transfer-exact-calibration-and-analytic-target}

For a declared Gaussian measure of literal mass seven and variance one,
psi=u\^{}2+1, q=2,n=2,N=3, the calculation gives

\begin{Verbatim}[breaklines=true,breakanywhere=true,fontsize=\small]
(omega_0,...,omega_4)=(7,7,14,42,168),
vartheta_1=22, vartheta_2=86/3,
Qhat_1=u, Qhat_2=u^2-11/3,
Delta_1=112, Delta_2=700/3,
G_3=diag(7/3,21/11), epsilon_3^2=400/99.
\end{Verbatim}

An independent quotient-Gram calculation with A=(1/2)I+i M\_u gives the
same result. The nonreal u-roots are algebraic fixtures, not zeros of
xi; the nonzero allowance explicitly tests that the method does not
manufacture purity. Repeated nodes, complex quartets, asymmetric atomic
tilt derivatives and mass changes are covered by the expanded checks.

For the actual quartet retain

\begin{Verbatim}[breaklines=true,breakanywhere=true,fontsize=\small]
q_k=[1+k(m-1)](k+1)^2,
L_(h,k)=2 delta [1+k(m-1)](k+1) floor((k+1)^2/4).
\end{Verbatim}

The source exterior/filtration result gives
L\_(h,k)\textless=epsilon\_(h,k,N), and L\_(h,k)/(k
q\_k)\textgreater=delta/2. Thus epsilon=o(k q\_k), with fixed h and
admitted N\textgreater=q\_k, is sufficient for the intended
contradiction. The new representation attaches that exact target to TWO
consecutive source-to-relation transfer coefficients and ONE phase
solve. It has NOT proved that asymptotic bound.

No new Lean execution or arithmetic interval certificate is claimed. All
new mathematical claims have the written proofs above at the retained
source inputs. The original arithmetic function, unit, action, theta
relations, residue pairing and trace remain unchanged.

\subsection[\hspace{0.4em}References and
checks]{References and
checks}\label{confluent-transfer-references-and-checks}

Classical background: C. Krattenthaler (2021), \textquotesingle A
determinant identity for moments of orthogonal polynomials that implies
Uvarov's formula for the orthogonal polynomials of rationally related
densities\textquotesingle, arXiv:2103.03969; NIST DLMF Section 18.2. The
proofs above supply the particular confluent and arithmetic interfaces
rather than attributing them wholesale to those references.

The expanded conversation delivery includes the complete LaTeX/HTML
reader and twenty exact finite regression methods; both normal and
optimized runs passed with matching records, and both deliberate-failure
controls failed. The preceding Toda archive\textquotesingle s 35
manifest entries verified and its 24-method checker passed fresh normal
and optimized runs. These are stated verification scopes, not new Lean
certificates. The public core checker in this directory independently
tests the principal raw-jet, transfer, norm, determinant, volume and
phase interfaces.

\endgroup
